\scanpage{62}
b) L'application $\Sigma^{<}(2)\xrightarrow{(p_1\circ s',\unsure{p_1b'|p_1\circ b'})}\Sigma\times\Sigma$
se factorise par $\Sigma^{<}(1)\lhook\joinrel\xrightarrow{(p_1,\uline{b})}\Sigma\times\Sigma$
en un morphisme
\[
p_2:\Sigma^{<}(2)-\Sigma^{<}(1)
\]
qui rende le carré
\diagram{\begin{tikzcd}[
  column sep={40mm,between origins},
  row sep={16.5mm,between origins},
  nodes={shape=rectangle,anchor=center,inner sep=0pt,outer sep=0pt},
  arrows={line width=.4pt,>=stealth,shorten <=0pt,shorten >=0pt},
  execute at end picture={
    \begin{scope}[shift={(s62a.center)},x=1cm,y=1cm,line width=.4pt]
    % UH1--UH4: the four separate upper hook fragments.
    \draw[line width=.4pt] (1.12,.22)--(1.06,.20);
    \draw[line width=.4pt] (.98,.15)--(.94,.12);
    \draw[line width=.4pt] (.94,.05)--(.98,.01);
    \draw[line width=.4pt] (1.05,-.02)--(1.12,-.02);
    % LH1--LH2: the two separate lower hook fragments.
    \draw (2.18,-1.40) .. controls (2.10,-1.39) and (2.08,-1.43) .. (2.06,-1.47);
    \draw[line width=.4pt] (2.07,-1.54)--(2.09,-1.58);
    % LS1--LS2: the lower shaft retains the gap from x=2.55 to x=2.62.
    \draw (2.21,-1.62)--(2.55,-1.62);
    \draw (2.62,-1.62)--(3.28,-1.62);
    % LR: the second, lower trace of the elongated open loop.
    \draw (3.24,-1.62) .. controls (3.58,-1.77) and (4.00,-1.73) .. (4.18,-1.62);
    % All four labels keep the accepted positions and underline scopes.
    \node at (2.35,.22) {$\uline{b}'$};
    \node at (.93,-.68) {$p_2$};
    \node at (4.75,-.65) {$p_1$};
    \node at (2.75,-1.32) {$\uline{b}$};
    \end{scope}
  }
]
|[alias=s62a]| \Sigma^{\unsure{<|*}}(2)
  \arrow[r,no head,to path={
    ([xshift=11.7mm,yshift=-.4mm]s62a.center)--
    ([xshift=34.5mm,yshift=-.4mm]s62a.center)
    .. controls ([xshift=35mm,yshift=-.4mm]s62a.center)
    and ([xshift=35mm,yshift=-.1mm]s62a.center)
    .. ([xshift=34.8mm,yshift=.5mm]s62a.center)
  }]
  \arrow[d,->,to path={
    ([xshift=2.5mm,yshift=-4.5mm]s62a.center)--
    ([xshift=10.3mm,yshift=-12.9mm]s62a.center)
  }]
&
|[alias=s62b]| \Sigma^{<}(1)
  \arrow[d,->,to path={
    ([xshift=42mm,yshift=-3.8mm]s62a.center)--
    ([xshift=48.7mm,yshift=-12.5mm]s62a.center)
  }]
\\
|[alias=s62c,xshift=15mm]| \Sigma^{<}(1)
  \arrow[r,->,to path={
    ([xshift=32.8mm,yshift=-16.2mm]s62a.center)
    .. controls ([xshift=36.5mm,yshift=-15.9mm]s62a.center)
    and ([xshift=40.2mm,yshift=-16mm]s62a.center)
    .. ([xshift=43.2mm,yshift=-16.2mm]s62a.center)
  }]
&
|[alias=s62d,xshift=10mm]| \Sigma
\end{tikzcd}}
\uline{Cartésien}, et $p_1:\Sigma^{<}(1)-\Sigma$ étant transversal
à \deleted{\unsure{l'imm}}ersion-but $\uline{b}:\Sigma^{\unsure{<|*}}(1)-\Sigma$.
\margin{\uline{NB} (Si $\Sigma$ est séparé, donc $\Sigma\times\Sigma$ séparé
sur $\Sigma$, et $\Sigma^{<}(1)\overset{\uline{b}}{-}\Sigma$
propre, il en résulte
que $\Sigma^{<}(1)-\Sigma\times\Sigma$ est
propre donc un plongement
\uline{fermé}.}
\margin{cette condition de transversalité est
peut être conséquence des autres .-}

De façon équivalente, une telle structure
est donnée par un espace topologique \inserted{séparé} $\Sigma$,
muni d'une relation d'ordre $x\preceq y$
\deleted{dont le graphe est}\inserted{satisfaisant les conditions suivantes}

% Two unassigned dark points beside the first cancelled a) line.
\margin{\figure{\begin{tikzpicture}[x=1mm,y=1mm]
\fill (0,0) circle (.16mm);
\fill (.5,-.8) circle (.16mm);
\end{tikzpicture}}}
\deleted{a) le graphe \deleted{$\unsure{\Sigma^{<}(1)|\Sigma^{*}(1)}$} $\Sigma(1)$ de cette relation est
fermé dans $\Sigma\times\Sigma$, et $\Sigma(1)\overset{\uline{b}}{-}\Sigma$ est
une immersion\unsure{.}
\deleted{$\unsure{\Sigma}$}
b) Si $\Sigma^{*}(1)$ est le graphe de la
relation d'ordre \inserted{stricte}, $\Sigma^{*}(1)\xrightarrow{\uline{b}}\Sigma$ est une
immersion \inserted{bordante} propre \deleted{\unsure{(donc la partie de}}
a) Si $x\in\Sigma$, $\Sigma_{<x}$ \inserted{$(=\Sigma^{\unsure{<|*}}(1)_{\unsure{x}})$}\unsure{,}
\unsure{\deleted{\unsure{est} \unsure{un|une}}|\unsure{est} \unsure{un|une}}
\unsure{ens|est} (fini) \uline{tot}$ {}^t$. \uline{ordonné}.
}
\margin{\handwritten{61}}
